% =========================
% Cas pratique PCT
% =========================

\subsection{Cas pratique}

%--------------------------
\begin{frame}{Cas pratique a)}

\begin{block}{quand?}
Délai de priorité : 12 mois (\pctart{8}.2.a)) + Art. 4.C.1 CUP.
\end{block}

En l'espèce, une demande nationale a été déposée le 17 mai 2023.\newline

\textbf{Conclusion: dépôt PCT sous priorité avant le 17/05/2024.}

\end{frame}

%--------------------------
\begin{frame}{Cas pratique a)}

\begin{block}{où?}
La demande internationale doit être déposée auprès d’un \textbf{office récepteur compétent},
c’est-à-dire :

\begin{itemize}
  \item un office national de l’État PCT où le demandeur est \textbf{domicilié}
  (\pctrule{19}.1(a)(i)) ;

  \item un office national de l’État PCT dont le demandeur est \textbf{ressortissant}
  (\pctrule{19}.1(a)(ii)) ;

  \item le Bureau international agissant en tant qu’office récepteur (RO/IB)
  (\pctrule{19}.1(a)(iii)) ;

  \item une organisation intergouvernementale (ex. OEB, ARIPO) ou un office national
  agissant pour l’État où un demandeur est domicilié ou ressortissant
  (\pctrule{19}.1(b)).
\end{itemize}
\end{block}

En l'espèce, A-FR est une société française. \newline

\textbf{Conclusion : RO = INPI, OEB, IB}

NB: pas d'obligation de dépôt RO/FR car pas première demande.

\end{frame}

\begin{frame}{Cas pratique a)}

\begin{block}{comment?}
\begin{itemize}
    \item déclaration de priorité: \pctrule{4}.10.a) + \pctrule{4}.1.b.i) (formulaire)
    \item demande internationale: \pctart{3}.2 (DI) + \pctart{4}.1 (requête) + \pctart{4}.1.d) (signatures).
\end{itemize}
\end{block}

\textbf{Conclusion: déposer une demande internationale comprenant:}
\begin{itemize}
    \item une requête comprenant une pétition PCT, désignations,
    nom \& renseignements du déposant et mandataire, titre invention,
    nom \& renseignements de l'inventeur;
    \item la requête est signée par déposant et mandataire;
    \item la copie de la demande FR.
\end{itemize}
NB: De préférence, requête sur formulaire PCT/RO/101.
\end{frame}

%--------------------------
\begin{frame}{Cas pratique b)}

\begin{block}{}
\begin{itemize}
  \item ajouter déposant après dépôt : \pctrule{92bis}.
  \item publication anticipée \pctart{21}.2.b).
\end{itemize}
\end{block}

En l'espèce, B-IE souhaite être rendue publique dès que possible.\newline

\textbf{Conclusion: ajouter co-déposant B-IE et
avancer la publication.}\newline
NB: éventuellement payer une taxe spéciale
si RRI pas disponible (\pctrule{48}.4) mais budget raisonnable.
\end{frame}

%--------------------------
\begin{frame}{Cas pratique c)}

\begin{block}{}
entrée en phrase nationale: \pctart{22}
\end{block}

En l'espèce, un produit suspect apparaît aux Pays-Bas.
Le niveau de développement de l'état de l'art est
nettement inférieur donc invention brevetable.
Pas de comparaison sur l'objet litigieux, éviter
une modification des revendications.\newline

\textbf{Conclusion : entrée en phase anticipée DO/EP.}\newline
NB: Etape suivante logique: valider brevet EP en EP/NL.
\end{frame}

%--------------------------
\begin{frame}{Cas pratique d)}

\begin{block}{}
\begin{itemize}
    \item entrée en phase anticipée \pctart{23}.3 + \pctart{40}.2
    \item PPH possible après ISR (pratique).
\end{itemize}
\end{block}

En l'espèce, CA-P veut accélérer la délivrance en ZA, US et MA.\newline

\textbf{Conclusion: ISR par EP => entrée en phase PPH(IP5) CN et US => AP Accord bilatéral CN-AP.}\newline
NB: PCT4ever

\end{frame}
